% !PRINTED 253
\weberpage{253}
Die Gruppe $E$ enth\"alt einen Theiler vom Index $p^n + 1$, der aus allen Substitutionen der Form
\[ \begin{pmatrix} \alpha, & \beta \\ 0, & \delta \end{pmatrix} \]
besteht, worin $\beta$ jede beliebige, $\alpha$ jede von Null verschiedene Zahl in $\mathfrak{C}$ bedeutet, und $\delta = \alpha^{-1}$ ist.

Diesem Theiler entspricht eine der Gruppe $E$ isomorphe Permutationsgruppe von $p^n + 1$ Ziffern, die man so bilden kann:

Der lineare Ausdruck
\weq{12}{\eta = \dfrac{\alpha\xi + \beta}{\gamma\xi + \delta} \quad \text{oder} \quad \xi = \dfrac{\delta\eta - \beta}{-\gamma\eta + \alpha}}
giebt f\"ur jede Zahl $\xi$ aus $\mathfrak{C}$ eine entsprechende Zahl $\eta$, ausgenommen, wenn $\gamma\xi + \delta = 0$ ist. Ebenso giebt es zu jedem $\eta$ ein bestimmtes $\xi$, ausgenommen f\"ur $\gamma\eta - \alpha = 0$. Um diese Ausnahme zu beseitigen, gen\"ugt es, das System $\mathfrak{C}$ durch Hinzuf\"ugung \gesperrt{eines} Elementes, das wir \glqq Unendlich\grqq{} nennen und mit $\infty$ bezeichnen, zu erweitern. Mit diesem Zeichen wird nach den folgenden Regeln gerechnet, worin $\alpha$ irgend ein Element des erweiterten Systemes $E$ bedeutet:
% [sic] Weber druckt gewoehnliches E sowohl hier als auch unten in "Systemes E", obwohl beide Stellen den erweiterten Koerper \mathfrak{C} meinen; originalgetreu bewahrt. Gewoehnliches E bezeichnet sonst die Congruenzgruppe.
\[
\begin{array}{l@{\quad}l}
\alpha + \infty = \alpha - \infty = \infty, & \text{ausser f\"ur } \alpha = \infty \\
\alpha \cdot \infty = \infty & \text{\textquotedbl}\ \ \text{\textquotedbl}\ \ \alpha = 0 \\
\alpha : 0 = \infty & \text{\textquotedbl}\ \ \text{\textquotedbl}\ \ \alpha = 0 \\
\alpha : \infty = 0 & \text{\textquotedbl}\ \ \text{\textquotedbl}\ \ \alpha = \infty.
\end{array}
\]
Dann f\"uhrt das Ergebniss jeder Rechnung mit den vier Species immer zu einer bestimmten Zahl des Systems $E$ und nur die Verbindungen $\infty \pm \infty, 0 . \infty, 0 : 0, \infty : \infty$ bleiben ohne Bedeutung.

Nach (12) entspricht dann der Zahl $\xi = -\delta : \gamma$ die Zahl $\eta = \infty$, und der Zahl $\xi = \infty$ die Zahl $\eta = \alpha : \gamma$, und jeder Substitution $A$ in $E$ entspricht eine bestimmte Permutation der $p^n + 1$ Elemente des erweiterten Systemes $\mathfrak{C}$. Man sieht leicht, dass nur die identische Substitution alle diese Elemente unge\"andert l\"asst, und dass folglich auch zwei verschiedene Substitutionen aus $E$ immer verschiedene Permutationen hervorrufen. Der Isomorphismus ist also einstufig.
